\documentclass[10pt,a4paper]{amsart}
\usepackage[margin=19mm]{geometry}
\usepackage{amsmath,amssymb,mathtools}
\usepackage[scr=rsfs,cal=euler]{mathalfa}
\usepackage{xfrac}
\usepackage[unicode,hidelinks,bookmarks=false]{hyperref}
\newcommand{\ie}{i.e.\ }
\newcommand{\cf}{cf.\ }
\newcommand{\resp}{respectively\ }
\newcommand{\Subsection}{\subsection}
\providecommand{\nobreakdash}{\nobreak\hbox{-}}
\setlength{\emergencystretch}{2em}
\multlinegap0pt
\usepackage{bm}
\usepackage{bbm}
\usepackage{dsfont}
\usepackage{xspace}
\usepackage{cancel}
\usepackage{mathtools}
\usepackage{amsthm}
\makeatletter
\@ifundefined{theorem}{\newtheorem{theorem}{Theorem}}{}
\makeatother
\title[Enumeration of intersection graphs of $x$-monotone curves]{Enumeration of intersection graphs of $x$-monotone curves}
\author{Jacob Fox, Janos Pach, Andrew Suk}
\date{Source: arXiv:2405.20547v2 (CC BY 4.0)}
\begin{document}
\maketitle

\noindent\emph{Source and licence.} Enumeration of intersection graphs of $x$-monotone curves. Version arXiv:2405.20547v2 (CC BY 4.0), distributed under the Creative Commons Attribution 4.0 International licence, as recorded in the arXiv licence metadata for this exact version and shown on the abstract page https://arxiv.org/abs/2405.20547v2. Full rights notice: licenses/arxiv-2405.20547-CC-BY-4.0.txt.\par\smallskip

\noindent\textit{Editorial scope.} Counted unit: the Theorem whose source label is \texttt{vc} in the article (source0.tex lines 300--302, source label \texttt{vc}), delivered together with the article's own complete original proof of it (source0.tex lines 305--323, one \texttt{proof} environment of 19 lines). No mathematics is written, repaired or re-derived: statement and proof are the article's own text line for line. Required context retained uncounted: hfromh' (lines 295--296). Every definition, display and label of those lines is retained unchanged. Omitted from this package and not redistributed: 1--294, 297--299, 303--304, 324--747. Nothing inside a retained excerpt is omitted or altered: every retained body is delivered exactly as the article's lines are written, so no omission receipt of any kind accompanies this package. Citations: the retained excerpts carry no citation command at all, so no bibliography entry of the article is transcribed and no reference is redistributed. Reference adaptation: none was needed and none was made; every reference inside the delivered unit resolves inside it, so no unresolved reference or citation is produced. Printed slips kept literal: no printed word, symbol, hypothesis or number of the counted statement or of its proof was altered. Layout and class: the article's own class is \texttt{article} with packages including amssymb, amsmath, amsthm, graphicx, geometry; the wrapper is the lane's pinned \texttt{amsart} editorial wrapper at 10pt on A4, which restates the theorem-like environments the retained text needs and adds the article's own macro definitions that the retained text needs (none), with extra wrapper packages (bm, bbm, dsfont, xspace, cancel, mathtools). The delivered numbering is the unit's own. Assets: no retained span contains a figure environment, a graphics-inclusion command, a PDF-page inclusion, a table or a TikZ picture; no image file is copied into this package, the package declares no asset dependency and no image file is redistributed with it. Licence: version arXiv:2405.20547v2 of this article is distributed under the Creative Commons Attribution 4.0 International licence, as recorded in the arXiv licence metadata for this exact version and shown on the abstract page https://arxiv.org/abs/2405.20547v2. A full rights notice is delivered at licenses/arxiv-2405.20547-CC-BY-4.0.txt. Characters: every retained excerpt is pure ASCII, so the delivered engine, XeLaTeX with its default Latin Modern text font, reports no missing character.
Let $h_{c,d}(m,n)$ denote the number of multiset systems $\mathcal{F}$ consisting of $m$ subsets of $[n]$ such that $\pi_{\mathcal{F}}(z) \leq cz^d$.   Let $h'_{c,d}(m,n)$ denote the number of set systems $\mathcal{F}$ of $m$ subsets of $[n]$ such that $\pi_{\mathcal{F}}(z) \leq cz^d$.  Clearly, $h'_{c,d}(m,n) \leq h_{c,d}(m,n)$.  It follows that $h'_{c,d}(m,n)=0$ if $m > cn^d$. Further, we can relate the two as follows. If we pick a multiset system consisting of $m$ sets whose primal shatter function is at most $cz^d$, then by throwing out repeated sets, we get a set system on the same ground set consisting of $m' \leq m$ sets. We then have to fill out these $m'$ sets to $m$ sets with repeats, including each set at least once. We thus have \begin{equation}\label{hfromh'} h_{c,d}(m,n)=\sum_{m' \leq m}h'_{c,d}(m',n){m-1 \choose m'-1}.
 \end{equation}

\begin{theorem}\label{vc}
Let $c,d \geq 2$ be fixed and $n,m \geq 2$. Then the number $h_{c,d}(m,n)$ of multiset systems $\mathcal{F}$ of $m$ subsets of $[n]$ such that $\pi_{\mathcal{F}}(z) \leq cz^d$,  satisfies $$h_{c,d}(m,n) = 2^{O(m^{1-1/d}n\log m)}.$$ Furthermore, if $m > cn^d$, then $$h_{c,d}(m,n)=2^{O(n^d\log m)}.$$
 \end{theorem}

 \begin{proof} 
Consider a linear ordering of the subsets of $[n]$. Let $\mathcal{F}$ be a multiset system of $m$ subsets of $[n]$. Let $S_1$ be the first set in $\mathcal{F}$ by the linear ordering. We will order the sets in $\mathcal{F}$ as $S_1,S_2,\ldots,S_m$ as follows. After picking $S_1,\ldots,S_{i-1}$, let $\delta_i=\max_{S \in \mathcal{F} \setminus \{S_1,\ldots,S_{i-1}\}} \min_{1 \leq j \leq i-1} d(S,S_j)$, and $S_i$ be a set $S$ that obtains the maximum, and $j_i$ be a $j$ that obtains the minimum $d(S_i,S_j)$. By our choice of the sets, 
the minimum of $d(S_a,S_b)$ over all $1 \leq a < b \leq i$ is $d(S_{j_i},S_i)$. By the Haussler packing lemma, we thus have $i = O((n/\delta_i)^d)$, or equivalently, $\delta_i = O(i^{-1/d}n)$. 

We now upper bound the number of choices of $\mathcal{F}$. There are at most $2^n$ choices of $S_1$. Each $j_i$ is a positive integer at most $i-1$, so there are at most $(m-1)! \leq m^m$ choices of $j_2,\ldots,j_m$. Having picked out this sequence of $j_i$'s, and having picked $S_1,\ldots,S_{i-1}$, we know $S_i$ must have symmetric difference at most $t_i:=O(i^{-1/d}n)$ from $S_{j_i}$. Thus, given this information, the number of choices for $S_i$ is at most $O(n^{t_i})$. Therefore, we obtain that the number of choices of $\mathcal{F}$ is at most 
\begin{eqnarray*} 2^n m^m \prod_{i=2}^m O(n^{t_i}) & \leq &  2^n m^m \prod_{i=2}^m (O(i^{1/d}))^{O(i^{-1/d}n)}\\ &= & 2^n m^m 2^{n\sum_{i=2}^m O(i^{-1/d}\log i)} \\ & = & 2^n m^m 2^{O(m^{1-1/d} n \log m)}.\end{eqnarray*} 
Note that the $2^n$ factor is at most the last factor. Hence, we obtain that the count is at most 
$m^m2^{O(m^{1-1/d} n \log m)}$. If $m \leq cn^d$, then the last factor is largest and this gives the desired bound. 

So we may assume we are in the case $m > cn^d$. In this case, by Equation (\ref{hfromh'}), the fact that $h'_{c,d}(m',n)=0$ for $m' > cn^d$ and $h'_{c,d}(m',n) \leq h_{c,d}(cn^d,n)$, we get $$h_{c,d}(m,n) \leq h_{c,d}(cn^d,n)\sum_{m' \leq cn^d}{m-1 \choose m'-1} = 2^{O(n^d \log m)}.$$
Notice that in this case, the first bound still holds, as $m^{1-1/d}n \log m \geq n^d \log m$. 







 \end{proof}

\end{document}
